%% ============================================================
%% RESUMES TRI-LANGUE -- une seule page, derniere du rapport
%% FR / EN / AR -- chaque section <= 4 lignes + mots-cles
%% ============================================================
\cleardoublepage
\thispagestyle{empty}
\pagestyle{empty}

\begingroup
\setlength{\parskip}{0pt}
\setlength{\parindent}{0pt}
\renewcommand{\baselinestretch}{1.1}
\footnotesize

% ---------- Français ----------
\begin{center}
  \textbf{\normalsize Résumé}
\end{center}
\noindent\textbf{Mots-clés :} IoT, sécurité, TLS~1.3, mTLS, PKI, HashiCorp Vault, MQTT, Suricata, Wazuh, ML, GNS3.

\smallskip
\noindent Ce projet, mené avec \TT{}, conçoit et valide une architecture de sécurisation de bout en bout d'un
écosystème IoT simulé sous GNS3. La solution associe une PKI HashiCorp Vault à deux niveaux, un broker MQTT
en TLS~1.3/mTLS obligatoire, un IDS Suricata et un SIEM Wazuh. Un pipeline ML (IsolationForest + RandomForest)
détecte les anomalies~: F1\,=\,0,994 et ROC-AUC\,=\,0,999 pour RF~. surcoût mTLS de l'ordre de \(+8\)\,\% sur le débit
applicatif et \(+3{,}5\)\,ms de latence, sans impact opérationnel notable.

\vspace{4pt}\noindent\rule{\linewidth}{0.4pt}\vspace{4pt}

% ---------- English ----------
\begin{center}
  \textbf{\normalsize Abstract}
\end{center}
\noindent\textbf{Keywords:} IoT, security, TLS~1.3, mTLS, PKI, HashiCorp Vault, MQTT, Suricata, Wazuh, ML, GNS3.

\smallskip
\noindent This project, carried out with Tunisie Telecom, designs and validates an end-to-end security architecture for a
simulated IoT ecosystem under GNS3. The solution combines a two-level HashiCorp Vault PKI, a Mosquitto MQTT
broker enforcing TLS~1.3/mTLS, a Suricata IDS and a Wazuh SIEM. An ML pipeline (IsolationForest + RandomForest)
performs anomaly detection: F1\,=\,0.994, ROC-AUC\,=\,0.999 (RF). mTLS overhead is about \(+8\)\,\% on application
throughput and \(+3.5\)\,ms latency, with no noticeable operational impact.

\vspace{4pt}\noindent\rule{\linewidth}{0.4pt}\vspace{4pt}

% ---------- Arabic ----------
\begin{center}
  \textbf{\normalsize \textarabic{ملخّص}}
\end{center}
\begin{Arabic}
\noindent\textbf{الكلمات المفتاحيّة:}
إنترنت الأشياء، أمن سيبراني، TLS~1.3، mTLS، PKI، HashiCorp Vault، MQTT، Suricata، Wazuh، تعلّم آليّ، GNS3.

\smallskip
\noindent صمّم هذا المشروع، المُنجَز مع \emph{اتّصالات تونس}، بنيةً لتأمين الاتّصالات من طرف إلى طرف داخل منظومة
إنترنت الأشياء المُحاكاة عبر GNS3. تجمع الحلّ بين بنية مفاتيح عموميّة (HashiCorp Vault)، ووسيط MQTT بـ TLS~1.3/mTLS،
ونظام Suricata، ومنصّة Wazuh SIEM. يُحقّق نموذج RandomForest في سلسلة التعلّم الآليّ
F1\,=\,0,994 وROC-AUC\,=\,0,999؛ ولا تتجاوز كلفة mTLS \(+8\)\,\% على الإنتاجيّة و\(+3{,}5\) ميلّي ثانية على زمن الاستجابة.
\end{Arabic}

\endgroup
